% Part: normal-modal-logic
% Chapter: syntax-and-semantics
% Section: relational-models

\documentclass[../../../include/open-logic-section]{subfiles}

\begin{document}

\olfileid{nml}{syn}{rel}

\olsection{تعلقی نمونے}

نارمل امکانی منطقاں دی معنیات دا بنیادی خیال \emph{تعلقی نمونہ}
اے۔ ایہ جہاناں دے اک سیٹ توں بندا اے۔ جہان دو رکنی «رسائی دے
تعلق» نال جُڑے نیں، تے اک قدر بندی دسدی اے کہ کیہڑے
!!{propositional variable}s کیہڑے جہاناں وچ «سچ» منّے جاندے نیں۔

\begin{defn}
  بنیادی امکانی زبان دا \emph{نمونہ} اک ترتیب وار تِکڑی
  $\mModel{M}
  = \tuple{W, R, V}$ اے، جتھے
  \begin{enumerate}
  \item $W$ «جہاناں» دا غیر خالی سیٹ اے،
  \item $R$، ~$W$ اُتے دو رکنی رسائی دا تعلق اے، تے
  \item $V$ اوہ فنکشن اے جیہڑا ہر !!{propositional
    variable}~$p$ نوں ممکن جہاناں دا سیٹ $V(p)$ دیندا اے۔
  \end{enumerate}
  جدوں $Rww'$ قائم ہووے، اسیں آکھدے آں کہ $w'$ تک
  \emph{رسائی}~$w$ توں اے۔ جدوں $w \in V(p)$، اسیں آکھدے آں
  کہ $p$،~$w$ \emph{وچ سچ} اے۔
\end{defn}

تعلقی معنیات دا وڈا فائدہ ایہ اے کہ نمونیاں نوں سادے خاکیاں
نال ظاہر کر سکدے آں، جیویں \olref{fig:simple} وچ۔ جہان گرہاں
نال ظاہر ہوندے نیں۔ جہان~$w'$ تک~$w$ توں رسائی اُدوں تے صرف
اُدوں اے جدوں~$w$ توں~$w'$ ول تیر ہووے۔ نالے اسیں گرہ،
یعنی جہان، اُتے~$\mTrue{p}$ لکھدے آں جدوں $w \in V(p)$،
تے ہور صورت وچ~\mFalse{p}۔ \olref{fig:simple} اوہ نمونہ
دکھاؤندا اے جیہدے وچ $W = \{w_1, w_2, w_3\}$،
$R = \{\tuple{w_1, w_2},\tuple{w_1,
  w_3}\}$، $V(p) = \{w_1, w_2\}$ تے $V(q) = \{w_2\}$۔

\begin{figure}
  \begin{center}
    \begin{tikzpicture}[modal]
      \node[world] (w1) [label={[align=right]right:\mTrue{p}\\ \mFalse{q}}]{$w_1$}; 
      \node[world] (w2) [label={[align=right]right:\mTrue{p}\\ \mTrue{q}},
        above right=of w1]{$w_2$}; 
      \node[world] (w3) [label={[align=right]right:\mFalse{p}\\ \mFalse{q}},
        below right=of w1] {$w_3$};
      \draw[->] (w1) to (w2);
      \draw[->] (w1) to (w3);
    \end{tikzpicture}
  \end{center}
  \caption{اک سادا نمونہ۔}
  \ollabel{fig:simple}
\end{figure}

\end{document}
